\section{Monomial pullback and face collision}\label{sec:pb1}

This section begins the extension beyond the current M0--M8 formalisation.
Its theorems are \evidence{PROVED} by the displayed finite-fibre arguments and
checked on released exact examples, but are not yet Lean theorems.

Let \(F(x)=\sum_{\alpha\in\N^n}A_\alpha x^\alpha\) be a formal series with
coefficients in a finite-dimensional vector space.  Fix an exponent matrix
\(\Lambda\in\N^{m\times n}\) with no zero column and non-zero leading constants
\(c_i\).  Define
\[
x_i=c_i y^{\lambda_i},\qquad
\lambda_i\in\N^m.
\]

\begin{theorem}[Fibre aggregation]\label{thm:aggregation}
The pulled coefficient at child exponent \(\beta\) is
\begin{equation}\label{eq:aggregation}
B_\beta(c)=\sum_{\Lambda\alpha=\beta}A_\alpha c^\alpha.
\end{equation}
The sum is finite.  Hence the actual child support is contained in, but may be
strictly smaller than, the naive image \(\Lambda(\supp F)\).
\end{theorem}
\begin{proof}
Substitution gives
\(A_\alpha c^\alpha y^{\Lambda\alpha}\).  Group equal child monomials to obtain
\cref{eq:aggregation}.  If \(\Lambda\alpha=\beta\), then for every column choose
a positive entry; this bounds each \(\alpha_i\) by a coordinate of \(\beta\),
so the fibre is finite.  Its sum can cancel.
\end{proof}

For a positive child weight \(\mu\), set \(w=\Lambda^{\mathsf T}\mu\) and let
\(\phi^{\rm gr}\) denote substitution by leading monomials.
This valuation transport is adjacent to tropical elimination
\parencite{SturmfelsTevelev2008}; the coefficient-fibre cancellation test is
kept explicit because equations can retain information lost by set-theoretic
support alone \parencite{Giansiracusa2016}.

\begin{theorem}[Candidate face and equality criterion]\label{thm:pullbackface}
If \(\nu_w(F)\) is the parent valuation, then
\[
\nu_\mu(\phi F)\ge\nu_w(F).
\]
The candidate component at that weight is
\begin{equation}\label{eq:candidateface}
\phi^{\rm gr}(\inw_w F).
\end{equation}
Equality holds if and only if \cref{eq:candidateface} is non-zero.  A literal
identity \(\inw_\mu(\phi F)=\phi^{\rm gr}(\inw_wF)\) without the non-vanishing
hypothesis is \evidence{FALSIFIED}.
\end{theorem}
\begin{proof}
Every parent monomial has child weight
\(\mu\cdot\Lambda\alpha=w\cdot\alpha\), so no pulled term can have smaller
weight.  Aggregating all parent terms of minimal weight gives
\cref{eq:candidateface}.  If it survives, it is the child initial form.  If it
cancels, the actual valuation is larger and the displayed right-hand side is
zero, whereas the child initial form, when the pullback is non-zero, occurs at a
later weight.
\end{proof}

\begin{figure}[ht]
\centering\begin{tikzpicture}[>=Latex,node distance=10mm and 18mm,
  box/.style={draw,rounded corners,minimum width=23mm,minimum height=10mm,align=center,font=\small},
  arr/.style={->,thick}]
\node[box,fill=blue!8] (a) {$A_\alpha x^\alpha$};
\node[box,fill=blue!8,below=of a] (ap) {$A_{\alpha'}x^{\alpha'}$};
\node[box,fill=orange!12,right=of $(a)!0.5!(ap)$] (b) {$B_\beta y^\beta$\\$=(A_\alpha c^\alpha+A_{\alpha'}c^{\alpha'})y^\beta$};
\node[box,fill=green!10,right=of b] (decision) {survives\\or cancels};
\draw[arr] (a)--node[above,font=\scriptsize] {$\Lambda\alpha=\beta$} (b);
\draw[arr] (ap)--node[below,font=\scriptsize] {$\Lambda\alpha'=\beta$} (b);
\draw[arr] (b)--node[above,font=\scriptsize] {test sum} (decision);
\end{tikzpicture}

\caption{Two parent exponents can enter the same child fibre.  Their aggregate,
not their separate presence, decides the child support.}
\label{fig:pullback}
\end{figure}

Under successive monomial maps with exponent matrices \(\Lambda\) then
\(\Omega\), exponents compose as \(\Omega\Lambda\); leading constants compose
as \(c_i d^{\lambda_i}\).  After a cancellation, however, the next survivor can
depend on the full germ and on unit jets, not merely on the first exponent
matrix.  No universal ``next face'' rule from the cancelled face alone is
claimed.
